\documentclass[12pt]{article}
\usepackage[T1]{fontenc}
\usepackage[utf8]{inputenc}
\usepackage{comment}
\usepackage{makecell}
\usepackage{graphicx} 
\usepackage{siunitx}
\usepackage{float}
\usepackage{url}

% 1-inch margins
\usepackage[margin=1in]{geometry} 

% Double spacing
\usepackage{setspace}
\doublespacing

% Use proper Times New Roman -- requires compiler XeLaTeX
%\usepackage{fontspec}
%\setmainfont{Times New Roman}

% set Times New Roman-style font but stay with pdfLatex compiler 
\usepackage{newtxtext}
\usepackage{newtxmath}

% For basic formatting
\usepackage{graphicx}
\usepackage{amsmath, amssymb}

% Optional: hyperlinks in references or text
%\usepackage[colorlinks=true, linkcolor=blue, parencitecolor=blue, urlcolor=blue]{hyperref}

% Page number in header
\usepackage{fancyhdr}
\pagestyle{fancy}
\fancyhf{}
\rhead{\thepage}
\renewcommand{\headrulewidth}{0pt}

% Paragraph formatting per APA
\setlength{\parindent}{0.5in}
\setlength{\parskip}{0pt}

% citations - APA 7
%\usepackage[backend=biber,style=apa]{biblatex}
\usepackage[style=apa,sorting=nyt,hyperref=true]{biblatex} % referencing. nyt = name, year, title
\addbibresource{bib.bib}

\begin{comment}
MAIN FOCUS FOR REWRITE:
* RESTRUCTURE INTRO FOR BETTER FLOW (outline below)
* CONSISTENT PHRASING PROBLEM-SKIPPING/HINT-REQUESTS/QM-USAGE
    - issue with this is that previous studies conceptualize this differently, so not sure if e.g. calling their measure 'problem-skipping' is correct.
    - look over paper to check how phrasing is latent-variable or observed-variable focused (accuracy vs. ability, problem-skipping vs. effort avoidance)
* ELABORATE META-COGNITION SECTION (explained below)
* CLARIFY THE INTERPRETATION OF THE CORRELATION BETWEEN THETA1 AND THETA2, BETA1 AND BETA2.

SMALL IMPROVEMENTS:
* Is connection / bridge between irt/bkt clear, and is it needed 
* Update referencing style 
* Fix fig. 8 (currently not shown in clockwise order)

OUTLINE OF INTRODUCTION
1. ADVANTAGES of unidimensional models and their contribution to educational measurement and practice
    - Invariant, interpretable measurement 
    - Scalable adaptive testing and practice 
    - Actionable diagnostics for teachers 
    - Efficient, motivated learning experiences
2. CRITICISMS of unidimensional modeling approaches 
    - not a criticism on unidimensional modeling for cogntive ability per se, but an explanation of how
    meta-cognitive states (induced by individual differences & platform design) can conflate ability estimates. 
3. THEORETICAL SECTION: Bridge into zero matters argument - one meaningful explanation to zero-state could be metacognition
4. EMPIRICAL LIT REVIEW: 
    - Findings of association problem-skipping and accuracy under unidimensional measurement models 
    - Measurement implications: MNAR, bias + empirical papers to support
    - Previous work in multidimensional modeling - offline contexts
5. MULTIDIMENSIONAL MODEL AS SOLUTION
    - explain model (equations)
6. SCIENTIFIC RELEVANCE
    - given 'zero matters' in online education...
    - given statistical issues with unidimensional models if underlying ability is not unidimensional...
    - we extend the use of IRTree from offline to online contexts (where more factors influence meta-cognition!)
    - implications 

ZERO-MATTERS ARGUMENT
* PRACTICAL PERSPECTIVE: There is a design issue. If apps offer a neutral response option, should they assume that choosing this response is neglible?
* STATISTICAL PERSPECTIVE: There is a modeling issue. Unidimensional models contain bias in measurement if different abilities exist. 
* THEORETICAL PERSPECTIVE: There is a variability issue. Individual differences in use of meta-cognitive strategies and risk-taking behaviors (and possibly others) result in different propensiities for using the netural response option --- given this variability, a unidimensional model will unequally and unfairly estimate ability. 
--> can also follow this reasoning / order of arguments in discussion 

-- Abstract -- 
core arguments
* unidimensional models are very powerful for estimating learner proficiency 
* they are also the foundation of CAT and adaptive learning systems
* ALSs provide a powerful, low-stakes practice environment which benefits student learning
* their low-stakes nature require extra design features which help learners stay engaged and motivated, such as skipping a problem (= non-responding)
* if non-response options are introduced, this can enable several behavioral pathways for a learner on a given item 
* choice of pathway depends on a learner's meta-cognitive state
* problem: unidimensional models do not account for meta-cognitive states. 
* because learners differ in their meta-cogntive dispositions, and response options may elicit different behaviors, the measurement of learner proficiency may be conflated with the learners 'meta-cognitive ability'
* we model an IRTree to separate cognitive ability from meta-cognitive ability (in this paper, in the form of problem-skipping). 

Outline
* Adaptive learning systems rely on accurate and fair assessment of student abilities to ensure optimal learning conditions.
* Importantly, nonresponse is often a design feature and bias gets amplified.
* if non-response options are introduced, this can enable several behavioral pathways for a learner on a given item 
* choice of pathway depends on a learner's meta-cognitive state
* problem: unidimensional models do not account for meta-cognitive states. 
* because learners differ in their meta-cogntive dispositions, and response options may elicit different behaviors, the measurement of learner proficiency may be conflated with the learners 'meta-cognitive ability'
* Method \& Result: we model an IRTree to separate cognitive ability from meta-cognitive ability (in this paper, in the form of problem-skipping). 
* Conclusion: From a measurement perspective, this requires a balance between powerful unidimensional approaches versus more complex models that can also capture individual differences in non-response processes. We give suggestions.

Drafts
%[EM:IP] In education, unidimensional measurement models have proven to be exceptionally powerful. By balancing psychometric rigor with practical usability, they underpin many large-scale assessments, high-stakes decisions, and adaptive testing and practice algorithms. However, this balance entails a trade-off that can have real consequences, particularly in adaptive learning settings where measurement and instructional decisions occur in rapid, continuous cycles. In this context, we show that meta-cognitive abilities can play a crucial role and should not be ignored, as they may amplify biased measurement. Specifically, we examine how problem-skipping behavior systematically affects measurement and the subsequent decisions based on it.

%Large-scale personalized learning systems rely on accurate and fair student ability estimations, yet common unidimensional models often neglect meta-cognitive processes. As such, they can provide biased estimates that put a limit on accuracy and fairness. This study introduces a multidimensional approach using the IRTree framework, which models problem-skipping as a distinct latent trait alongside accuracy. Problem-skipping is a common meta-cognitive manifestation that can reflect self-regulation and strategic decision-making. Our results demonstrate that problem-skipping reflects meaningful individual differences, with patterns of problem-skipping varying in their impact on performance. Notably, while frequent problem-skipping was associated with lower accuracy, it also revealed adaptive strategies for some learners. This suggests that problem-skipping is signal, not noise. To ensure that learning analytic tools are supported by accurate and fair measurement models, we suggest several ways to account for problem-skipping when estimating student ability.
%In reality, however, individuals differ in their propensity to use meta-cognitive strategies, thus conflating estimates of cognitive ability. 
%Adaptive learning systems rely on accurate and fair assessment of student abilities to ensure optimal learning conditions. On the other hand, to keep learners engaged, they allow for different types of non-responding. From a measurement perspective, this requires a balance between powerful unidimensional approaches versus more complex models that can also capture individual differences in non-response processes. This balance entails a trade-off that can have real consequences, particularly in adaptive learning settings where measurement and instructional decisions occur in rapid, continuous cycles. In this context, we show that meta-cognitive abilities can play a crucial role and should not be ignored, as they may amplify biased measurement. Specifically, using a large dataset from a K12 adaptive learning system, we examine how problem-skipping behavior systematically affects measurement and the subsequent decisions based on it.

%Unidimensional measurement models are a powerful tool for estimating learner proficiency in educational settings. Thanks to their psychometric rigor and practical usability, they have become foundational for large-scale adaptive learning systems. However, to maintain engagement and motivation during practice, many adaptive learning systems utilize non-response options. While beneficial for the learner experience, these options may come at a cost, b

\end{comment}


\begin{document} 

%\title{From noise to signal: Separating problem-skipping from accuracy in student ability estimation}
%\title{A Problem that Shouldn't be Skipped: Problem Skipping Limits the Accuracy of Ability Estimates in Online Learning}

% [BJET] Accurately tracking learner abilities with non-cognitive variation: Accounting for differences in non-response
% Non-response ontspoort : Accounting for metacognitive .. in large-scale adaptive systems

% What Skipping Tells Us: Non-Response, Metacognition, and Fairness in Adaptive Educational Technology
% Meta-cognitive dynamics in adaptive learning
% A problem that shouldnt be skipped: modeling non-response in adaptive educational technology 
% ...: Separating Ability From Non-Response in Adaptive Educational Technology

% The Problem with Skipping Problems: Separating Cognitive From Metacognitive Dynamics in Adaptive Educational Technology
% What's the Problem with Problem-Skipping? Separating Cognitive From Metacognitive Dynamics in Adaptive Learning Systems

%% my current top title contenders :-) 
% When Learners Skip a Problem: Separating Cognitive From Metacognitive Dynamics in Adaptive Learning Systems
% When Learners Skip a Problem: Disentangling Ability and Non-response in Adaptive Learning Systems 


\title{When Learners Skip a Problem: Disentangling Ability and Non-Response in Adaptive Learning Systems}
\date{}
\maketitle


\begin{abstract}
% word count EM:IP: cannot find exact limit, but there are no recent papers with >150 words. 
% word count BJET: also cannot find, but published abstracts are considerably longer, 200-300 words. 
% currently 200 words
%Adaptive learning systems rely on fast, fair and accurate assessments of student abilities to ensure an optimal matching between learners and problems. For this purpose, unidimensional measurement models are highly effective and widely adopted. However, as opposed to traditional measurement settings, adaptive learning systems typically allow for different types of non-responding to keep learners engaged. This may amplify meta-cognitive dynamics, harm unidimensionality and consequently reduce accuracy and fairness. Additionally, the continuous updates of ability estimates that drive such adaptive learning systems may further amplify these issues. In this study, we use a tree-based item response model to separate ability from non-responding by analyzing the responses of over \num{3500} learners from a large-scale K-12 online learning platform. We demonstrate that ability and non-responding are moderately correlated, but distinguishable latent traits, confirming that non-responding should not be ignored, as it amplifies biased measurement of ability. Ensuring optimal learning conditions requires a balance between powerful unidimensional approaches versus more complex models that can also capture individual differences in metacognitive processes. We suggest several ways to deal with this trade-off. 

% Importantly, introducing these design features enables several behavioral pathways for a learner on a given item, the choice of which depends on their meta-cognitive state at a given point in time. This gives rise to a measurement problem: under the unidimensional model, individual differences in cognitive- and non-response processes are captured by the same latent trait, while in reality, response tendencies might differ across learners.
% These findings suggest that non-response options carry meaningful information about learner behavior and should not be collapsed into cognitive ability alone

% outline (from alexander)
%1. Problem-skipping and non-response are a problem, particularly in adaptive edtech.
%2. This is in part a design consideration, for motivational reasons.
%3. This has wanted and unwanted metacognitive effects.
%4. This is a problem for unidimensional models that power adaptivity.
%5. Implications with respect to fairness/accuracy and possibly for design considerations.

%% rough draft
%- while learning, observable behavior from students may not only reflect cognitive abilities but also meta-cognitive regulation. 
%- this is especially the case in adaptive educational technology platforms, which often allow for different types of non-responding to keep learners engaged. 
%- while beneficial from a motivational perspective, metacogntivie dynamics induced by non-response patterns may harm the assumption of unidimensionality which most educational measurement models rely on. 
%- this has consequences for accuracy and fairness in adaptive systems. 
%- additionally, the continuous updates of ability estimates that drive adaptive learning systems may further amplify these issues. 
%- in this study, we use tree-based item response model to separate ability from non-responding by analyzing the responses of over 3500 learners from a large-scale K-12 online learning platform. we demonstrate that ability and non-responding are moderately correlated, but distinguishable latent traits, confirming that non-responding should not be ignored when estimating learner ability. 
%- ensuring optimal learning conditions requires a balance between powerful unidimensional approaches versus more complex models that can also capture individual differences in metacognitive processes.
%- educational technology platforms need to consider how to account for metacognitive dynamics while ensuring robust estimates of cognitive abilities. we suggest several ways to deal with this trade-off.

%% refined
Student learning behavior reflects not only cognitive ability, but also metacognitive regulation---how learners monitor and manage their own learning. This is particularly relevant in adaptive educational technology platforms, which often allow for different types of non-responding, such as skipping a problem. Commonly, these options are an intentional design choice to sustain engagement. While such features can be beneficial for motivation, the non-response patterns it induces may also introduce metacognitive dynamics that threaten the assumption of unidimensionality underlying most educational measurement models. This has consequences for the accuracy and fairness of ability estimates, and these risks may be amplified by the continuous updating of estimates that underpin adaptive systems. %This has consequences for the accuracy and fairness of ability estimates, which threaten to be further compounded by the continuous updating of estimates that drives adaptivity. 
In this study, we apply a tree-based item response model to disentangle cognitive ability from non-responding behavior, using data from over \num{3500} learners from a large-scale K-12 adaptive learning platform. We find that ability and non-responding are moderately correlated, but distinguishable latent traits, confirming that non-responding should not be ignored when estimating learner ability. These results highlight a core tension in adaptive educational technology design: powerful unidimensional models offer practical efficiency, but more complex models may be needed to capture individual differences in metacognitive processes. We discuss implications for both measurement practice and platform design, and propose several ways to navigate this trade-off. 
 
\end{abstract}

\paragraph{Keywords}
adaptive learning, learner performance modeling, problem-skipping, item response tree, K-12 education, fairness

\section{Introduction}
% plan-- BJET: write a general introduction here, set up the premise of the paper
% - explain adaptive learning platforms
% - explain how they rely on ability estimates (require continuous measure of ability)
% - introduce design choice that creates the problem 

% [§0.1 - BJET]
Adaptive learning technologies are increasingly used to support personalized instruction in digital learning environments. By continuously adjusting the difficulty of tasks to match a learner's current level of understanding, such systems aim to provide practice that is neither too easy nor too difficult. To accomplish this, adaptive platforms rely on statistical models that estimate a learner's current proficiency based on their observed performance during practice. Statistical models continuously update these ability estimates, which simultaneously guide the selection of subsequent tasks. Consequently, the accuracy of these models directly affects how effectively instruction can be personalized.

% [§0.2 - BJET]
% Design choice -- many systems offer alternative response options, introduces alt behavioral pathways etc etc. 
% End this par with problem statement (most can be moved from §1.3 EM:IP)
At the same time, many digital learning environments offer response options beyond simply answering a question. Learners may request hints, skip items, or indicate that they do not know the answer. Such options are often included to support productive learning strategies, reduce frustration, or maintain engagement during practice. However, they also introduce multiple behavioral pathways through which learners can interact with a task. When learners choose to skip a problem or request help rather than attempt a response, their behavior may reflect metacognitive regulation, strategic effort allocation, or disengagement rather than task proficiency alone. This raises an important question for adaptive learning systems: should such behaviors be interpreted separately from learner ability? 

% [§1.1 - BJET]
\subsection{Unidimensional measurement models}
Adaptive learning systems estimate learner proficiency from response data using statistical measurement models such as item response theory (IRT) and Bayesian knowledge tracing (BKT). Although these frameworks differ in their formulation, they share a common assumption: observable performance is driven by a single latent dimension representing task proficiency. Such unidimensional measurement models have fundamentally transformed educational measurement in practice. They offer a parsimonious and interpretable framework for representing student proficiency~\parencite{rasch1993probabilistic, hambleton1991fundamentals, embretson2013item, de2013theory} and under adequate model fit, they enable the comparison of individuals' ability regardless of which specific items from a calibration pool that these students respond to~\parencite{rasch1993probabilistic, wright1982rating}. This ability has been ground-breaking for large-scale assessment~\parencite{lord2012applications, van2010elements}. With these tools, educators and test constructors can identify poorly functioning items while maintaining robustness against, for example, missing data~\parencite{mislevy1996missing, de2013theory}.  

% [§1.1 - EM:IP]
%Unidimensional measurement models, such as the Rasch model and closely related item response theory (IRT) formulations, have fundamentally transformed educational measurement in practice. They offer a parsimonious and interpretable framework for representing student proficiency~\parencite{rasch1993probabilistic, hambleton1991fundamentals, embretson2013item, de2013theory} and under adequate model fit, they enable the comparison of individuals' ability regardless of which specific items from a calibration pool that these students respond to~\parencite{rasch1993probabilistic, wright1982rating}. This ability has been ground-breaking for large-scale assessment~\parencite{lord2012applications, van2010elements}. With these tools, educators and test constructors can identify poorly functioning items while maintaining robustness against, for example, missing data~\parencite{mislevy1996missing, de2013theory}.  

% [§1.2 - shared]
These advantageous properties have made unidimensional models the backbone of computerized adaptive testing and practice systems, which rely on a single proficiency scale to dynamically match learners to appropriately targeted items~\parencite{weiss1982improving, van2010elements, magis2017computerized}. In instructional contexts, this common metric enables teachers to monitor progress over time, manage mastery, and interpret student performance relative to item difficulties~\parencite{corbett1995knowledge, maris2012speed, deonovic2018learning, saric-grgic2024twenty-five}. Likewise, at the curriculum level, psychometric modeling supports evidence-centered refinement of learning objectives~\parencite{mislevy1996missing, glaser2001knowing}. For students, adaptive learning systems can successfully align task difficulty with estimated proficiency, and their use is associated with improved learning efficiency and sustained engagement~\parencite{eggen2006optimal, ma2014intelligent, van2010elements}. 

% [§1.3 BJET -- needs to be rewritten to correspond with intro section]
Although adaptive learning systems built on these unidimensional approaches appear powerful for education, their advantages rest on a critical assumption: observable performance is driven by a single latent dimension. Under this assumption, all observable outcomes are treated as manifestations of the same underlying ability, effectively collapsing multiple learning processes into a single scale. If response behaviors are driven by processes other than proficiency, unidimensional models may conflate distinct sources of variability, potentially biasing both person and item estimates. For example, when learners have the option to skip a question, one student may skip because they correctly judge that they do not know the answer, whereas another may skip simply to avoid the effort of responding. If the measurement model treats these behaviors as equivalent indicators of proficiency, differences in strategic behavior may become embedded in the estimated ability parameter. Over time, systematic differences in the use of the skip option could therefore bias ability estimates. In this paper, we examine this fundamental measurement question for adaptive learning systems: does non-responding represent noise around ability, or signal a distinct latent process that should be modeled separately? 

% [§1.3 EM:IP]
% Although adaptive learning systems built on these unidimensional approaches appear powerful for education, their advantages rest on a critical assumption: observable performance is driven by a single latent dimension. Learning platforms often offer options beyond simply answering a question correctly or incorrectly, such as requesting hints or skipping items. These features are necessary for engagement and motivation, but also introduce additional behavioral pathways that may reflect metacognitive regulation, strategic effort allocation, or disengagement rather than true task proficiency. Unidimensional models treat all observable outcomes as manifestations of the same underlying ability, thus statistically collapsing all learning processes into a single scale. When such behaviors are driven by processes other than proficiency, unidimensional models may conflate distinct sources of variability, potentially biasing both person and item estimates. For example, when learners have the option to skip a question, one student may skip because they correctly judge that they do not know the answer, whereas another may skip simply to avoid the effort of responding. If the measurement model treats these behaviors as equivalent indicators of proficiency, differences in strategic behavior may become embedded in the estimated ability parameter. Over time, systematic differences in the use of the skip option could therefore bias ability estimates. In this paper, we examine this fundamental measurement question for adaptive learning systems: does non-responding represent noise around ability, or signal a distinct latent process that should be modeled separately? 

\subsection{A metacognitive interpretation of non-responding} % interpretation? view? perspective? explanation? 
%\subsection{Psychometrics meets metacognition: Why non-responding matters} 
% set up the link with metacognition
When faced with an online educational exercise, for example answering the multiplication problem ``3 $\times$ 4'', a student must decide whether and how to respond. This decision involves monitoring one's knowledge state (\textit{e.g.,} "Do I know the answer?"), evaluating potential outcomes (\textit{e.g.,} "What happens if I answer incorrectly?"), and regulating effort allocation (\textit{e.g.,} "Is this worth attempting?"). Such processes fall under the umbrella of metacognition, which involves the monitoring and regulation of cognition during learning~\parencite{winne2017cognition}. Importantly, this regulation may result in the decision not to answer the exercise at all for some learners while others would attempt to answer the items. While metacognition has been shown to play a role for successful learning in offline~\parencite{rabinowitz2017nteraction, fleur2021metacognition, colantonio2024predicting} and online~\parencite{zhao2020time, kim2009not, fidan2024emotions, shields2024active} settings, it is clearly not reducible to proficiency alone. Thus, non-responding is not merely an absence of performance --- it is an observable trace of monitoring and regulation. In this sense, "zero matters"~\parencite[p.~2]{van2025statistical}. 

%* THEORETICAL PERSPECTIVE: There is a variability issue. Individual differences in use of meta-cognitive strategies and risk-taking behaviors (and possibly others) result in different propensiities for using the netural response option --- given this variability, a unidimensional model will unequally and unfairly estimate ability. 
The extent to which meta-cognitive tools are used differs across learners, as a function of individual differences in, for example, prior knowledge~\parencite{taub2014can}, motivation~\parencite{sungur2007modeling, karlen2016differences}, and anxiety~\parencite{matthews1999metacognition, wells1995metacognition}. Differential patterns in non-responding may also arise from avoidance strategies linked to affective states such as anxiety or boredom, which influence engagement and self-regulation through their effects on perceived control and motivation~\parencite{pekrun2013control}. Compounding this issue, students do not always recognize when help is needed, leading to systematic differences in help-seeking and response behavior~\parencite{aleven2000limitations}. Taken together, these findings suggest that the decision not to respond may reflect stable individual differences in metacognitive regulation rather than proficiency alone. %If such differences are ignored, systematic patterns of problem-skipping may become conflated with ability, potentially biasing the estimation of student proficiency.

%* PRACTICAL PERSPECTIVE: There is a design issue. If apps offer a neutral response option, should they assume that choosing this response is neglible?
% [§2.3 BJET Practical perspective]
Non-responding is not only driven by individual responding behavior but it is also structurally embedded in the design of many adaptive learning systems~\parencite{corbett1995knowledge, aleven2000limitations, vanlehn2011relative, roll2015understanding}. When a platform deliberately offers a non-response option, it actively shapes the behavioral space in which learners engage. Such a design choice carries implicit measurement assumptions about what type of information about the learner that the non-response behavior holds (for instance, when learners' ability estimates are not affected by problem-skipping, there is an implicit assumption that skipping a problem carries no information about ability). Yet, as the individual differences reviewed above suggest, this assumption is unlikely to hold uniformly across learners. This ambiguity is further amplified by the adaptive nature of the system itself, which continuously updates ability estimates and uses them to select subsequent items~\parencite{van2010elements}. If non-responding is conflated with proficiency in this feedback loop, the ability estimate may systematically misguide learners by selecting items in the wrong difficulty range, which in turn generates further distorted response patterns. 
The design of the learning system thus plays a vital role in shaping metacognitive behaviors---these are not only driven by individual differences, but also conditioned and amplified over time by design. % reformulate closing sentence: design and individual differences go hand in hand -- and are amplified by adaptive nature. 

% [§2.3 EM:IP Practical perspective]
%Non-responding is not only meaningful at the individual level; it is also structurally embedded in the design of many adaptive learning systems. In addition to selecting an answer, learners are often given alternative response options such as skipping an item, requesting hints, indicating uncertainty, or making partial attempts~\parencite{corbett1995knowledge, aleven2000limitations, vanlehn2011relative, roll2015understanding}. These features are usually introduced to support productive learning processes, for example by reducing unproductive guessing, scaffolding help-seeking, or allowing learners to regulate effort and difficulty. At the same time, the presence of such options introduces multiple behavioral pathways through which learners can interact with an item and in doing so, complicate the interpretation of non responses. This raises an important design and measurement question: when a learning platform offers a neutral response option, can the resulting non-response be treated as negligible, or as equivalent to learner proficiency?

%* STATISTICAL PERSPECTIVE: There is a modeling issue. Unidimensional models contain bias in measurement if different abilities exist. 
The handling of a non-response is also a statistical question. In many measurement models, non-responses are either ignored or treated as incorrect responses. Ignoring non-responses relies on the assumption that missingness is random, \textit{i.e.,} the response is missing-at-random (MAR). Under this assumption, the probability of a non-response does not depend on the unobserved response value, conditional on the observed data~\parencite{rubin1976inference}. Alternatively, treating a non-response as incorrect assumes that the processes leading to non-response and incorrect responses are governed by the same latent trait, such as proficiency. Given variability in the behavioral pathways induced by non-response options and learners’ individual differences in metacognitive regulation, it is likely that non-responding depends, at least partly, on processes other than the target latent trait. In such cases, the missingness mechanism becomes non-ignorable, \textit{i.e.,} missing-not-at-random (MNAR), and collapsing non-responses into a single latent trait estimate may introduce bias in measurement. Addressing this issue may therefore require modeling non-response behavior as a distinct latent process alongside proficiency~\parencite{little2019statistical}. Whether this distinction holds in the specific context of adaptive learning platforms---where non-response options are intentionally designed into the system---remains an open empirical question.

%Together, these theoretical, behavioral, and statistical considerations suggest that non-response behavior in adaptive learning environments may reflect a process distinct from task proficiency. Next, we review empirical evidence for this distinction.

\subsection{Previous literature}
Research on non-responding in online learning environments is limited, and existing studies provide a mixed picture of its relationship with proficiency. Results from a large-scale intervention study showed that limiting access to a problem-skipping option increased effortful practice and improved learning outcomes in an online mathematics platform \parencite{savi2018delaying, savi2023combating}. Similarly, hint requests have been found to correlate negatively with pre- and post-test mathematics scores and positively with the number of errors made during practice \parencite{norum2024student}. These findings suggest that the use of non-response options may be associated with task proficiency. 

At the same time, several studies highlight that the use of such options varies substantially across learners. For example, \textcite{norum2024student} found that students who were identified as needing the most support were also the least likely to request a hint. Eye-tracking evidence similarly indicates that students with lower ability devote less attention to available hints \parencite{conati2013understanding}. Other work has documented substantial individual differences in how learners process hints, although these differences were not related to learning gains \parencite{goldin2012learner, goldin2013hints}. In addition, studies aiming to detect disengaged or off-task behavior in intelligent tutoring systems have identified groups of students who misuse non-response options \parencite{baker2004off, baker2004detecting, baker2005performance, beck2004using}. Taken together, these findings suggest that problem-skipping behavior cannot be explained by student ability alone.

If non-response behavior reflects processes beyond proficiency, treating these responses as ignorable or equivalent to incorrect answers may distort measurement. Indeed, psychometric research has demonstrated that ignoring non-responses in IRT models can introduce bias in both person and item parameter estimates~\parencite{finch2008estimation, mislevy1996missing, pohl2014dealing, holman2005modelling, debeer2017modeling}. Given the reviewed literature which indicates that non-responding and learner proficiency likely are related, it is possible that non-responding distorts measurement of proficiency when it is ignored. To address this issue, several approaches have been proposed to explicitly model response processes alongside the substantive construct of interest. One such approach is the Item Response Tree model (IRTree; elaborated in Section \dots), which represents the response process as a hierarchical sequence of decisions \parencite{bockenholt2012modeling, de2012irtrees, thissen2013two}. In this framework, the observed response is generated by multiple latent processes corresponding to different decision stages. IRTree models have been successfully applied to capture differences in response styles in survey data \parencite{bockenholt2017response, thissen2013two}.

Applications of IRTree models in educational assessment have provided evidence that non-responding and proficiency may represent distinct processes. Simulation studies show that unidimensional models produce biased person and item parameters when responses are MNAR, whereas explicitly modeling the missingness mechanism with an IRTree model substantially reduces this bias \parencite{debeer2017modeling}. Empirical studies of PISA reading assessments in Argentina \parencite{debeer2017modeling} and Japan \parencite{okumura2014empirical}, as well as studies of reading and mathematics competence among German fifth-grade students \parencite{pohl2014dealing}, found that the propensity to skip an item was negatively related to performance and that IRTree models provided a better fit than standard IRT models. These results indicate that non-responses in educational assessments may reflect processes distinct from proficiency.

IRTree models have also been applied to online learning data to account for behavioral processes such as the speed–accuracy trade-off in ability estimation \parencite{coomans2016distinguishing, hofman2018fast}. However, these applications did not consider problem-skipping behavior. Consequently, most existing work on skipping has been conducted in traditional testing contexts. Whether similar multidimensional structures arise in adaptive learning environments---where skipping and help-seeking are embedded in platform design---remains largely unexplored. 

\subsection{Multidimensional solution}
IRTree models present one approach to modeling non-responses in an educational context. In specific, the model treat the interaction for a learner on a given item as a two-stage decision process: first, the item is either responded to or not, and second, \textit{given that the item was responded to}, the response is correct or incorrect. This process can be visualized as a tree-branching model with two nodes (Figure \ref{fig:tree}), and the response of person $p$ on item $i$ can be coded as two binary variables, $Y^{(1)}$ and $Y^{(2)}$, where 

\begin{equation}
Y^{(1)} =
\begin{cases} 
1 & \text{if the person attempts item } i \\
0 & \text{otherwise}
\end{cases}
\end{equation}

\begin{equation}
Y^{(2)} =
\begin{cases} 
1 & \text{if the person solves item } i \text{ correctly,} \\
0 & \text{otherwise}
\end{cases}
\label{eq:cases}
\end{equation}

The tree model is a multidimensional IRT model, where the probability of the response of person $p$ on item $i$ being in category $m$ can be formulated as: 

\begin{equation}
\pi(x_{pi} = m \mid \theta_p, \beta_i) = \prod_{n=1}^{N} \left[ \left( \frac{1}{1 + e^{-(\theta_p^{(n)} - \beta_i^{(n)})}} \right)^{t_{mn}} \right]^{q_{mn}}
\label{eq:tree}
\end{equation}

Each node in the tree is modeled with a 1PL model. Depending on where you are in the tree, the probability is computed as the product across all internal nodes involved in the respondent's path through the tree, leading to a specific response category. The terms $t_{mn}$ and $q_{mn}$ decide whether a node will contribute to the probability of category $m$, in our case question-mark, incorrect or correct. $t_{mn}$ determines whether an internal node $n$ is relevant for the probability calculation of response category $m$. It equals 1 if the node is relevant to category $m$ (\textit{i.e.}, contributes to the response), and 0 otherwise.
$q_{mn}$ determines whether the node $n$ is active for the probability calculation of category $m$). When $q_{mn} = 1$, the node will actively contribute to the probability of category $m$, and when $q_{mn} = 0$, the node will be ignored for that category.

For example, consider a student who does \textit{not} skip item $i$ ($Y^{(1)} = 1$), and then responds correctly ($Y^{(2)} = 1$): at node 1, their probability of attempting the item (not skipping) is modeled as a 1PL model (equation \ref{eq:Rasch}), and at node 2, their probability of correctly answering item $i$ follows a second 1PL model (equation \ref{eq:Rasch}). The final response depends on successfully moving through both nodes, so the overall probability of providing a correct answer, $P(Y^{(1)}=1, Y^{(2)}=1)$, is the product of these two probabilities. At both nodes, $t_{mn}$ and $q_{mn}$ are both equal to 1, as they are relevant and active for the final response category. Inversely, for a student who \textit{does} skip the problem (\textit{i.e.}, $Y^{(1)} = 0$), node 1 is relevant ($t_{m1} = 1; q_{m1} = 1$) but node 2 is not ($t_{m2} = 0; q_{m2} = 0$). Thus, the overall response probability is only the probability of skipping ($P(Y^{(1)} = 0)$), modeled by a 1PL model, and they do not proceed further in the tree. 

\begin{figure}[t]
        \centering
        \includegraphics[width=0.3\linewidth]{EDM-25/figures/tree.pdf}
        \caption{A binary decision tree (IRTree model) with 2 nodes reflecting the decision to respond or not ($Y_1$) and whether the response is correct or not ($Y_2$).}
        \label{fig:tree}
\end{figure}

The main advantage with this modeling approach is that the unidimensionality of the decision process can be formally tested~\parencite{de2012irtrees}. The two-node branching model depicted in Figure~\ref{fig:tree} gives rise to two latent traits on the person and item side respectively: 
\begin{enumerate}
    \item node 1: $\theta_p^1$ (individual propensity to skip an item) and $\beta_i^1$ (propensity for item to be skipped), and 
    \item node 2: $\theta_p^2$ (individual ability) and $\beta_i^2$ (item difficulty).
\end{enumerate}
Now, the correlation between $\theta^1$ and $\theta^2$, and $\beta^1$ and $\beta^2$ can be calculated. That is, how much do skipped responses depend on respondents or items, respectively? The unidimensional model assumes that these correlations are negligible, \textit{i.e.}, for individuals, problem-skipping and accuracy are the same, and the likelihood for an item to be skipped is equal to its estimated difficulty. This assumption is tested by comparing the unidimensional to the multidimensional (IRTree) model. 

\subsection{The current study}
To date, a hierarchical decision model, allowing problem-skipping and accuracy to be modeled as separate latent traits, has not been applied to large-scale adaptive learning data.
% reiterate 2 main points (theory & model)
Leveraging behavioral patterns in the form of skipped responses can give insight into learner meta-cognition and disengagement, and by providing a more accurate model of student ability can contribute to more fairness in adaptive learning systems. 

In this work, we estimate the relationship between problem-skipping and accuracy in a large-scale adaptive educational software designed for primary school students to practice arithmetic. We compare unidimensional (1PL) and multidimensional IRTree models with varying levels of parameter constraints. These models allow us to examine whether responses to items are driven by one latent trait (ability) or whether accuracy and problem-skipping represent distinct underlying processes:

\begin{enumerate}
    \item[\textbf{RQ1:}] Are problem-skipping and accuracy distinguishable processes in online learning?  
\end{enumerate}

Finding evidence for multidimensionality in ability assessment would have large implications for current educational measurement, which largely assume ability as a unidimensional construct. In addition, fitting the IRTree model to response data allows us to examine the latent correlations between nodes for items and students. This would help shed light on the currently sparse findings concerning the relationship between accuracy and problem-skipping: 

\begin{enumerate}
    \item[\textbf{RQ2:}] Given that problem-skipping and accuracy are distinct, how are they related? 
\end{enumerate}


%%%% not incorportated %%%%

% Par 4 Explanation IRT and connection to BKT. 
%Problem-skipping has previously been incorporated into ability assessment, largely within Bayesian Knowledge Tracing (BKT) frameworks. Since its introduction~\parencite{corbett1995knowledge}, where problem-skipping behavior is treated the same as an incorrect response, there have been several extensions which treat problem-skipping differently. For example, the partial-credit model~\parencite{wang2010representing, wang2013extending, ostrow2015improving}, gives some, but not full, credit for hint-seeking or multiple attempts to a problem. Other implementations include incorporating the history of help-seeking to improve the prediction accuracy of guesses and slips~\parencite{baker2008more} or detecting affective states~\parencite{corrigan2015dynamic}. A comprehensive overview of such proposed enhancements to the classical BKT model is provided by~\textcite{saric-grgic2024twenty-five}. 

%Another commonly used measurement model for assessing student ability is Item Response Theory (IRT). In this model, student ability is assessed based on the likelihood of a correct response, given specific properties of the current test item, and the student's current ability. In this paper, we use extensions of the 1-Parameter Logistic Model (1PL; also referred to as the Rasch Model~\parencite{rasch1993probabilistic}) to account for hierarchical decision processes. In the 1PL model, the probability that a student with ability $\theta_p$ answers item $i$ correctly is given by the function: 
%\begin{equation}
%  P(X_i = 1 \mid \theta_p, \beta_i) = \frac{1}{1 + e^{-(\theta_p - \beta_i)}}
%  \label{eq:Rasch}
%\end{equation}
%where $\beta_i$ is the difficulty parameter for item $i$, which represents at what ability level the student has a 50\% chance of responding to the item correctly. In BKT terms, the $\theta$ parameter (student ability) represents the probability of moving from an unlearned to a learned state, and the $\beta$ parameter (item difficulty) represents an item-specific probability of moving from a learned state back into an unlearned state. This bridge between IRT and BKT models of student ability has been formalized by~\parencite{deonovic2018learning}. Although the current paper deals with data obtained under an IRT measurement model, results will be relevant for online learning systems utilizing a BKT-framework. 

%Importantly, IRT and BKT models share the feature that ability is measured on a unidimensional scale. Both measurement models thus assume that problem-skipping is driven by the same latent trait as the accuracy of a response. Previous work has demonstrated that ignoring missing responses in IRT approaches can lead to bias on person- and item-parameter estimates~\parencite{finch2008estimation, mislevy1996missing, pohl2014dealing, holman2005modelling, debeer2017modeling}. Similarly, prediction accuracy was found to increase in BKT models when disengaged responses (including long sequences of skipped problems) were excluded from the ability estimation~\parencite{gorgun2022considering}. These findings support the notion that problem-skipping might be indicative of a separate meta-cognitive state which is not directly related to ability, and signal a need to test the assumption that problem-skipping and accuracy can be measured on the same dimension. 

%Failure to account for individual differences in response patterns is problematic because it may bias the validity of the trait measurement and other constructs to be estimated,  particularly when the response style is related to the trait to be measured. Therefore, a range of research is dedicated to reliably identifying and controlling for response styles to produce ability estimates which are as unbiased as possible. One established way is to model the response process as a hierarchical decision tree, referred to as an IRTree model~\parencite{de2012irtrees}, where the content and the response format of the item are both assumed to contribute to the observed response~\parencite{bockenholt2012modeling, de2012irtrees, thissen2013two}.

\section{Methods}
\subsection{Data}
\subsubsection{Online Learning Environment}
% ANONYMOUS VERSION
We use data from an adaptive learning platform, [X] developed for primary school students to study Math and Language \footnote{In the current version of this manuscript, the name of the learning platform, and the country in which it is active, is hidden for anonymity. In places where the final manuscript would refer to the platform, the name has been replaced with [X].}. The platform currently consists of 65 games divided between environments designed for Arithmetic practice (28 games), language learning (24 games), and for second-language English learning (13 games). It is currently being used by over \num{450000} children across \num{4300} schools. Depending on the difficulty of the task, the student can select one out of multiple answer options, or fill in the number themselves. They can skip the item by clicking a question mark. Then, the correct answer is shown and the student continues to the next item. While on the platform, children take care of a virtual garden, where each game represents a different learning domain, and the amount of practice is shown as the health of the plant. 

In this work, we use data from a game called Series (Figure~\ref{fig:series}), wherein a student is asked to fill in a missing number in a number sequence. This game is designed to test knowledge of number sequences and relations. Each item that a student completes is given a score, determined by an explicit scoring rule (Figure~\ref{fig:scoring}, Equation~\ref{eq:scoring}). This scoring rule is visualized to the student in the form of coins: coins disappear for every second until the deadline (20 seconds). If the item is answered correctly, the amount of coins remaining is awarded, whereas if the answer is incorrect, the remaining amount of coins is lost. If the student skips the item or does not answer in time, no change is made to their score. 

% NON-ANONYMOUS VERSION
%We use data from the adaptive online learning environment Prowise Learn (previously known in the literature, and referred to here as, Math Garden~\parencite{klinkenberg2011computer}). The learning platform was developed for Dutch primary school students to study Math and Language. The platform currently consists of 65 games divided between the learning environments Math Garden (for Arithmetic practice; 28 games), Language Sea (for Dutch language learning; 24 games), and Words and Birds (for English language learning; 13 games). It is currently being used by over \num{450000} children across \num{4300} schools. Depending on the difficulty of the task, the student can select one out of multiple answer options, or fill in the number themselves. They can skip the item by clicking a question mark. Then, the correct answer is shown and the student continues to the next item. While on the platform, children take care of a virtual garden, where each game represents a different learning domain, and the amount of practice is shown as the health of the plant. More information about the platform is available at \href{https://www.prowise.com/en/product/prowise-learn/}{the Prowise Learn Website}. Research on the adaptivity of the platform can be found in~\parencite{klinkenberg2011computer, brinkhuis2018learning, brinkhuis2020governing, hofman2018solution}. The platform has also been used for applied developmental and educational research, across topics such as skill development in mathematics~\parencite{ou2019modeling} and language~\parencite{kievit2019mutualistic}, math anxiety~\parencite{jansen2013influence, jansen2016affective, hilz2023relationship}, error-monitoring~\parencite{de2022post, de2021error}, effort~\parencite{savi2023combating} and problem-skipping~\parencite{savi2018delaying}. 

%In this work, we use data from the Math Garden game Series (Figure~\ref{fig:series}), wherein a student is asked to fill in a missing number in a number sequence. This game is designed to test knowledge of number sequences and relations. Each item that a student completes is given a score, determined by an explicit scoring rule (Figure~\ref{fig:scoring}, Equation~\ref{eq:scoring}). This scoring rule is visualized to the student in the form of coins: coins disappear for every second until the deadline (20 seconds). If the item is answered correctly, the amount of coins remaining is awarded, whereas if the answer is incorrect, the remaining amount of coins is lost. If the student skips the item or does not answer in time, no change is made to their score. 

Item difficulties and student abilities are estimated with the Elo Rating System (ERS). When entering a learning game for the first time, a student is given a provisional ability rating $\theta$ (around the mean of the existing students within the same grade level). For each item that the student responds to, $\theta$ is updated according to the weighted difference between the accuracy and the expected accuracy on the item. Simultaneously, the rating of the item ($\beta$) is updated: 

\begin{align}
    \hat{\theta}_p &= \theta_p + K_p(S_{ip} - E(S_{ip})), \\
    \hat{\beta}_i &= \beta_i + K_i(E(S_{ip}) - S_{ip}).
\end{align}
% wrap it in \begin{equation} to have it appear as one equation (but error messsage occurs)

\noindent $K$ is a scaling parameter which determines the weight that the difference between the observed and expected score will have on the new ability estimate (for more information about the K-factor, see~\textcite{vermeiren2024balancing}). The score of person $p$ following their response $x$ on item $i$ is dictated by the difference between their response time $t$ and the deadline $d$, scaled by a discrimination parameter $\alpha$:

\begin{equation}
    S_{ip} = (2x_{ip} - 1)(\alpha_id_i - \alpha_it_ip).
    \label{eq:scoring}
\end{equation}

\noindent This is the formalization of the explicit scoring rule (Figure~\ref{fig:scoring}). It is a so-called high speed, high stakes (HSHS) scoring rule~\parencite{maris2012speed}, which imposes a speed accuracy trade-off in response behavior. The discrimination parameter, $\alpha_i$, is equal to $1/d$. In this way, the model maintains a constant discrimination value across items, and thus the calculated expected score is equivalent to a 1PL IRT model: 

\begin{equation}
    E(S_{ip}) = a_i d_i \frac{e^{2 a_i d_i (\theta_p - \beta_i)} + 1}{e^{2 a_i d_i (\theta_p - \beta_i)} - 1} - \frac{1}{\theta_p - \beta_i}.
\end{equation}

% ANONYMOUS VERSION
\noindent The ERS adaptively matches children to items according to their desired difficulty level, which correspond to the probability of answering an item correctly. Difficulty levels are 60\% (hard) 75\% (medium), or 90\% (easy). Children play on the default medium setting unless they select otherwise. 

% NON-ANONYMOUS VERSION
%\noindent The ERS adaptively matches children to items according to their desired difficulty level, which correspond to the probability of answering an item correctly. Difficulty levels are 60\% (hard) 75\% (medium), or 90\% (easy). Children play on the default medium setting unless they select otherwise. For an in-depth explanation of how the ERS has been applied to Math Garden, see~\parencite{klinkenberg2011computer}.

\begin{figure}[htbp]
    \centering
    \begin{minipage}{0.5\textwidth}
        \centering
        \includegraphics[width=\linewidth]{EDM-25/figures/series.jpeg}
        \caption{A screenshot of the Series game. A student is asked to fill in the missing number in the sequence. They gain (correct answer) or lose (incorrect answer) coins corresponding to the remaining amount of time in seconds (bottom right). A student can click the question mark (middle right) to retrieve the answer to the item and move on to the next.}
        \label{fig:series}
    \end{minipage}%
    \hfill % Add space between the figures
    \begin{minipage}{0.5\textwidth}
        \centering
        \includegraphics[width=\linewidth]{EDM-25/figures/scoringrule.png}
        \caption{Visualization of the explicit scoring rule.}
        \label{fig:scoring}
    \end{minipage}
\end{figure}

\subsubsection{Inclusion criteria}
To obtain data suitable for the IRT fitting procedure we selected a specific subset of data from the Series game. First, we sampled data from the 2-year period between 01--09--2014 and 31--08--2016. In recent years, changes have been made to the explicit scoring rule and the accessibility of the question mark response in [X]. To ensure that these changes do not affect question mark behavior in our analyses, we chose to look at data before this period of time. Lastly, the adaptive nature of the online learning environment means that user ability parameters are continuously updated. This also implies that for each user, there are periods of stable (learning is constant) and unstable (denoted by jumps or shifts in learning) measurements of user ability. In order to control for possible effects of learning on question mark usage, we only included data from periods where a user's ability remained relatively stable. We defined a stable period as one where the user's ability rating stayed close to their average rating. Specifically, a stable period was when the user's ability rating did not fluctuate more than 1 rating point above or below their average rating. The average rating for each user was calculated by taking the 20\% trimmed mean of all their ability ratings, to prevent outliers from skewing the average. This criterion helps to minimize the influence of learning (or sudden jumps in ability) on question mark usage. Lastly, for proper identification of the models, we selected only users who had skipped a minimum of 10 responses within their data period. The length of each obtained period of learning data varied between users, with the mode around 1--2 weeks and the max up to 80 weeks.

On the item side, we selected items which had been played a minimum of \num{2000} times within the selected data collection period. The final dataset consisted of \num{3795} respondents and \num{386} items. 

\subsection{Analyses}
\label{section:analyses}
Comparing unidimensional and multidimensional linear response tree models allows to test whether one underlying trait gives rise to the response categories on a test item~\parencite{de2012irtrees}. Here, we compare a unidimensional (1PL) model to IRTree models with varying parameter constraints to investigate whether responses to items in the Series game rely solely on one latent person variable (ability) or whether accuracy is distinguishable from problem-skipping. Because IRTree models are forms of generalized linear mixed models, the software package \textit{lme4}~\parencite{bates2015package} is used to estimate them. In all models, response is a binary variable denoting a decision (left or right) in the response tree. Each response in the tree is modeled as a 1PL model. The response tree is formulated such that node 1 represents the propensity to skip a problem, with 1 coded as a question mark response and 0 coded as an attempt to solve the item. Node 2 represents the propensity to give an incorrect response (difficulty), with 1 coded as an incorrect response and 0 coded as a correct response. See Figure~\ref{fig:tree} and Equations~\ref{eq:cases} and~\ref{eq:tree} for the technical specifications of the IRTree models. As such, each model is multilevel logistic regression which predicts response as a function of users and items, with varying constraints on the item- and user- specific random slopes: 

\begin{enumerate}
    \item Fully estimated IRTree: A multidimensional linear response tree where response is predicted by a random node effect of items and a random node effect of users. This model captures item- and user-specific variation across nodes. 
    \item Item-constrained IRTree: A multidimensional linear response tree where response is predicted by a random intercept for items, and a random node effect of users, but without modeling item-specific variation in nodes. This model captures variability in user behavior across nodes, but not item thresholds. 
    \item User-constrained IRTree: A multidimensional linear response tree where response is predicted by a random intercept for users, and a random node effect of items, but without modeling user-specific variation in nodes. This model captures variability in item thresholds across nodes, but not user behavior. 
    \item Fully constrained IRTree: A unidimensional random item model for linear response trees, where both item- and user-level effects are modeled with random intercepts, but without accounting for variation across nodes. Item omission and correctness are collapsed into the same response process, making this conceptually equivalent to a 1PL model. 
\end{enumerate}

% Also include how we compare these models or is that redundant? 

\section{Results}
\subsection{Description of Problem-Skipping }
On average, children skipped approximately 6\% of items in the Series game. Our exploration of problem-skipping behavior revealed several trends. First, unsurprisingly, the usage of the question mark feature was used more by children playing on the hard level ($M = 0.086, SD = 0.281$), compared to easy ($M = 0.043, SD = 0.203$) and medium ($M = 0.050, SD = 0.218$) levels. Additionally, there was a small increase in the overall use of question mark responses across grades (0.061 (SD = 0.239) for Grade 3, 0.059 (SD = 0.237) for Grade 4, 0.062 (SD = 0.241) for Grade 5, 0.062 (SD = 0.241) for Grade 6, 0.063 (SD = 0.244) for Grade 7, and 0.067 (SD = 0.250) for Grade 8.). However, when looking at average question mark usage across grades and difficulty levels simultaneously, this relationship changed, with stable proportions of problem-skipping for children playing on the hard level, and a decreasing proportion of problem-skipping across grades for easy and medium levels (Figure~\ref{fig:grade_diff}). 

\begin{figure}[htbp]
      \centering
      \includegraphics[width=0.5\linewidth]{EDM-25/figures/grade_diff.png}
      \caption{Proportion of question mark responses in the Series game, separated by grade and difficulty level (easy = 90\% correct; medium = 75\% correct; hard = 60\% correct). Ribbons around the line denote the $95\%$ confidence interval. Grade refers to the Dutch grade levels; age ranges approximately between 4 and 12 years. QM = question mark. }
      \label{fig:grade_diff}
\end{figure}  

We also found differences with regards to at what point in time the question mark feature was used. Figure~\ref{fig:rt} displays the probability of each response type (correct, incorrect, question mark) across response times. Problem-skipping was much more likely to occur in the first five seconds upon the presentation of an item, and very unlikely to occur after five seconds, though with a small increase in probability upon nearing the response deadline at 20 seconds. This is in stark contrast to the probability of making a correct response, which quickly increased in the first five seconds and slowly decreased leading up to the deadline. The probability of making an incorrect response showed a small peak early in the response process, and slowly increased with increasing response times. This figure also demonstrates that within the response times that are most common (where the histogram bars are the highest), the probability of giving a question mark response is relatively low. 

Similarly, with increasing number of items played in the Series game, the probability of skipping the problem decreased (Figure~\ref{fig:session}). This indicates that as users gain more understanding of how to play the Series game, their problem-skipping decreases. This graph also demonstrates the adaptivity of the system: as user ability estimates get more certain, the probability of making an error and question mark response stabilizes. 

\begin{figure}[htbp]
    \centering
    \begin{minipage}{0.48\textwidth}
        \centering
        \includegraphics[width=\linewidth]{EDM-25/figures/rt.png}
        \caption{The probability of each response type (correct, error, question mark) across response times (seconds). Point size refers to the amount of observations at that data point. The vertical line denotes the deadline to respond, at 20 seconds. The histogram displays the average distribution of response times in the Series game.}
        \label{fig:rt}
    \end{minipage}%
    \hfill
    \begin{minipage}{0.48\textwidth}
        \centering
        \includegraphics[width=\linewidth]{EDM-25/figures/session.png}
        \caption{The probability of answering with a question mark or giving an incorrect response, over the first 100 items played in the Series game. Point size refers to the amount of observations at that data point.}
        \label{fig:session}
    \end{minipage}
\end{figure}



\subsection{Models of Problem-Skipping }
To test whether problem-skipping is distinguishable from accuracy in the Series game, we compared a fully estimated multidimensional response tree to a unidimensional response tree, and two versions of a multidimensional response tree where either random item estimates (item-constrained IRTree) or random user estimates (user-constrained IRTree) were constrained (full model specifications are highlighted in Section~\ref{section:analyses}). We compared these models on the basis of AIC and BIC model fit (Table~\ref{tab:model-comp}). The fully estimated IRTree model fit the data best, supporting the inclusion of a latent trait for problem-skipping in the estimation of student ability. 

The fully estimated IRTtree model had a fixed intercept of $-1.85$ ($p < .001$). This translates to a baseline probability of making an incorrect response of approximately $14.1\%$. Crucially, the correlation between the propensity to problem-skip (node 1) and give an incorrect response (node 2) was $0.41$: students that were more likely to problem-skip were moderately more likely to make erroneous responses. Although this correlation is positive, it is not 1 (which is the assumption under the unidimensional model). Similarly, the correlation between the item skipping threshold and item difficulty was $0.77$. That is, items that were more likely to be skipped were more likely to be estimated as difficult. The correlations between nodes for items and users is visualized in Figure~\ref{fig:cor}. 

\begin{table*}[t]
    \centering
    \small % Reduce font size for the table
    \begin{tabular}{p{1cm}p{3.5cm}ccccc}
        \hline
        \textbf{Model} & \textbf{Random Parameters} & \textbf{AIC} & \textbf{BIC} & \textbf{RMSE} & \textbf{cor($\theta^{(1)}, \theta^{(2)}$)} & \textbf{cor($\beta^{(1)}, \beta^{(2)}$)} \\
        \hline
        1 & $\theta_p, \theta_i, \beta_p, \beta_i$ & 1519007 & 1519095 & 0.3217 & 0.41 & 0.78 \\
        2  & $\theta_p, \beta_p, \theta_i = \beta_i$ & 1533720 & 1533783 & 0.3234 & 0.35 & - \\
        3  & $\theta_i, \beta_i, \theta_p = \beta_p$ & 1550148 & 1550211 & 0.3251 & - & 0.73 \\
        4 & $\theta_i = \beta_i, \theta_p = \beta_p$ & 1566326 & 1566377 & 0.3268 & - & - \\
        \hline
    \end{tabular}
    \caption{Model fit statistics and latent correlation estimates for the four fitted models. RMSE = root mean square error. RMSE values were computed via $k$-fold cross-validation, where 10\% of the data was held out in each fold, the model was fit on the remaining data, and RMSE on the model predictions was calculated on the held-out set. The reported RMSE is the average across all folds.}
    \label{tab:model-comp}
\end{table*}


\begin{figure}[t]
    \centering
    \includegraphics[width=0.75\linewidth]{EDM-25/figures/cor.pdf}
    \caption{Correlation between estimated nodes in the fully estimated IRTree model. The left graph displays the correlation between the item skipping threshold ($\beta^{(1)}$), and the item difficulty ($\beta^{(2)}$). The right-hand graph displays the correlation between the individual propensity to skip an item ($\theta^{(1)}$), and the individual propensity to answer incorrectly ($\theta^{(2)}$). For both graphs, The scatter plot denotes individual data points, and the regression line displays the best-fitting linear relationship between the nodes. Density plots denote the distribution of each random parameter.}
    \label{fig:cor}
\end{figure}

\subsection{Individual Patterns in Problem-Skipping }
\begin{figure*}[t]
    \centering
    % \includegraphics[width=1\linewidth]{user_interactions.pdf}
    \includegraphics[width=\textwidth,height=6cm]{EDM-25/figures/user_interactions.pdf}
    \caption{Eight individual users extracted based on their estimated propensity to skip and make an incorrect response (center graph). Individual interaction patterns are displayed in a clockwise fashion starting from the upper right corner (low accuracy, high skipping), to the upper left corner (low accuracy, low skipping) of the scatter plot. Each graph plots each each individual's response sequence to items in the Series game. The position on the y axis reflects the response time to the item. Colors denote the response type (correct, incorrect, question mark). While some users had longer response sequences in the extracted data (users C and E), the x axis is limited to showing the first 300 responses, to ease the visualization.
    }
    \label{fig:interactions}
\end{figure*}

In addition to indicating the average relationship between individual propensities to problem-skip and give an incorrect response, the scatter plot of user estimates in Figure~\ref{fig:cor} reveals the spread of data around this relationship. From this, it could be of particular interest to look closer at the play behavior of students that either conform to or deviate from the average pattern of problem-skipping and accuracy. Exploring individual patterns in problem-skipping could reveal further insights and generate new hypotheses about individual differences in interacting with online learning systems. To explore this, we randomly selected users from 4 sets of criteria: (1) Low accuracy and high skipping rate (upper right corner of the scatter plot); (2) High accuracy and high skipping rate (lower right corner of the scatter plot); (3) High accuracy and low skipping rate (lower left corner of the scatter plot); (4) Low accuracy and low skipping rate (upper left corner of the scatter plot). We visualized their interactions with items in the Series game by plotting their response type (correct, incorrect, question mark) for each item across time, against their response time (Figure~\ref{fig:interactions}). This allows us to look at patterns in response types across time as well as whether responses were fast or slow and how this fluctuates. 

This exploration revealed some interesting patterns. For example, because user ratings cannot improve with a question mark response (see Figure~\ref{fig:scoring}), users with a high ability should, in theory, not resort to using the question mark response at a high rate. When taking a closer look at these users (users C and D in Figure~\ref{fig:interactions}), they seem to be resorting to a different strategy compared to the other sampled users: long sequences of question-mark responding to avoid giving an answer, combined with sequences of attempting the item wherein the answer is mostly correct. In fact, the majority of users within this region of average skipping and accuracy showed such sequences, whereas the majority of users outside this region did not. 
 
\section{Discussion}
%% NOTES %%%%%%%
% I discuss the person-level correlation here, but the item-level correlation is not really discussed. A reviewer could question what the 0.77 item-level correlation implies? Perhaps the unidimensional model is a reasonable approximation for items? Is this relevant?

Accurate educational assessment demands accurate measurement models. In this study, we compared a unidimensional IRT model, which assumes that problem-skipping does not influence student ability differently from accuracy, to a multidimensional IRTree model, which explicitly accounts for differences in propensity to skip an item. We found evidence that problem-skipping and ability are distinguishable processes, and as such, should not be measured on the same scale. Additionally, we found a moderate, positive association between problem-skipping and accuracy, suggesting that students with a high propensity to skip problems are somewhat more likely to be inaccurate when they do give a response. However, the relative weakness of this correlation opposes the assumption of unidimensional measurement models that these are perfectly correlated. Instead, it implies that skipped responses carry information about both accuracy and other meta-cognitive or motivational traits which affect responding. Failure to account for this will introduce bias to measurement and limit the adaptivity of the system. 

Our findings extend previous work on response scale- and achievement test data---demonstrating that ignoring skipped items may lead to bias in the measurement of latent traits~\parencite{holman2005modelling, pohl2014dealing, huang2020mixture, okumura2014empirical, debeer2017modeling}---to online learning platforms. They also shed some light on the mixed evidence concerning the relationship between problem-skipping and accuracy, adding to the findings of~\textcite{norum2024student} that more problem-skipping is associated with more errors, and the body of literature indicating individual differences in the use of alternative response options~\parencite{norum2024student, goldin2012learner, goldin2013hints}. However, the evidence that problem-skipping and accuracy are distinct and only moderately related may also help to explain the discord in the literature. Specifically, the imperfect relationship suggests that the observed association between problem-skipping and accuracy may be influenced by contextual factors, such as the design of the learning environment or task difficulty, which give space for different manifestations of meta-cognitive strategies and motivational states such as frustration or boredom. Future work should aim to disentangle these contextual influences to better understand the role of problem-skipping in learning and assessment. 

Many adaptive learning platforms include a response option alternative to providing a direct answer to an item, such as skipping or requesting a hint. These findings therefore have widespread relevance for any adaptive learning system that aims to provide a fair estimate of student ability while allowing flexibility in how students engage with items. A basic recommendation for such systems is to, at a minimum, measure both latent traits once the data have been collected. The IRTree model has proven a feasible method for this kind of post-hoc analysis, offering insights into individual differences in both accuracy and problem-skipping. However, there are some limitations to the use of IRTree models: they require large amounts of data for reliable estimates, making them less suitable for students who interact with the learning platform less frequently. Moreover, in the way that it is currently implemented, this method reflects past behavior rather than providing real-time insights, a core feature in adaptive learning systems. In practice, the IRTree framework could be used to incorporate a second parameter into the measurement model for real-time tracking of problem-skipping. This would enable adaptive systems to monitor and act on these tendencies as they emerge, for example via teacher dashboards. While this method is still data-intensive, it would allow for dynamic tracking of both accuracy and problem-skipping, providing a fuller picture of the learning process. 

Despite their importance in revealing individual differences in how problem-skipping relates to accuracy, post-hoc and real-time models reveal little about where these differences stem from. This raises the question of whether different types of meta-cognitive strategies or states--such as boredom, frustration, gaming, or wheel-spinning--can be inferred from students' response data. Consider two students, one of whom quickly skips items in a very long sequence (students C and D, Figure \ref{fig:interactions}) to avoid effort, and another who is exerting effort into learning by attempting as many items as they can, but skips the item as soon as they are uncertain to avoid the risk of negative feedback. While both students would display high problem-skipping tendencies, they bode for different intervention styles and may therefore want to be separated. To aid in this, more information can be added to the tree model, such as a node to distinguish fast from slow responses~\parencite{coomans2016distinguishing, hofman2018fast}. An IRTree model can also be extended to model different levels of disengagement, by, for example, explicitly accounting for drop-out~\parencite{huang2020mixture}. 

Apart from merely measuring different cognitive and meta-cognitive states from observed response data, adaptive systems can be designed to make these strategies directly observable. One approach is to provide response options explicitly linked to meta-cognitive strategies, allowing systems to infer and adapt to the strategies students employ. The feasibility of such an approach would depend on the scale and flexibility of the learning system. For instance, in a system with an exploratory style of learning, mapping response options to every possible strategy may prove challenging. An alternative approach is to change the scoring rule in such a way that it constrains students' behaviors to a specific strategy. For example, in [X], the HSHS scoring rule~\parencite{maris2012speed} is implemented in such a way that students adhere to a speed-accuracy trade-off. In a similar manner, a separate scoring rule for problem-skipping can be implemented such that it has an appropriate trade-off with accuracy. A potential future research avenue is to investigate the applicability and validity of such approaches for detecting different problem-skipping strategies. 

Lastly, if problem-skipping is seen as undesirable behavior--such as indicating low effort or gaming the system--the design of the system or scoring rule can be changed in such a way that this behavior is decreased. Such an intervention has been successfully implemented, wherein the ability to use the problem-skipping option was delayed upon the presentation of each item~\parencite{savi2018delaying}. This led to an increase in the amount of items that students attempted, promoting effortful responding in the system. Such interventions highlight the potential of system design to nudge students toward more productive behaviors.

A final note to consider is that this model assumes that the student undergoes a sequential decision process, \textit{i.e.}, first deciding whether to respond to a problem or skip it, and only then deciding on the actual response (if not skipped). It is not certain that this is the true decision process of the student. In the future, it would be interesting to compare different modeling frameworks aimed at deciphering the response process of students. For example, Race Models~\parencite{heathcote2022winner} assume that there is competition between two or multiple parallel processes, and as enough evidence accumulates in support of one process, this process wins and determines the subsequent decision. In the context of problem-skipping, such a model may be suitable for detecting fine-grained processes such as individual differences in processing speed, especially in the context of a speed accuracy trade-off~\parencite{heathcote2022winner, evans2018modeling}. Alternative methods to model students' decision-making processes include Bayesian Hierarchical Models~\parencite{sawyer2018modeling, lee2014bayesian} and Markov Decision Processes~\parencite{lamar2018markov, puterman1994markov}. The rich, granular data generated by online learning systems offer an ideal opportunity to compare cognitive decision models, which usually require detailed and structured data collected in lab experiments. Nevertheless, the current results already indicate that there are structural individual differences which, regardless of how the decision procedure is modeled, needs to be accounted for.

%\vspace*{2em} % add if header hovers alone at end of page in final version
\section{Conclusion}
In a large-scale adaptive learning environment for arithmetic practice, we found that accuracy and problem-skipping are moderately related, but distinct traits. Under a measurement model which doesn't explicitly account for individual differences in problem-skipping, student ability estimates could be biased towards children who respond differently to problem-skipping options. To prevent this, online learning environments should consider estimating and reporting latent traits for both the tendency to problem-skip and give an accurate response, or intervene in the learning system such that strategies can be inferred or problem-skipping is reduced. We have shown the feasibility of applying the IRTree framework~\parencite{de2012irtrees} to model problem-skipping in an adaptive learning platform. This approach has the potential to be extended to model other meta-cognitive strategies, or model problem-skipping in domains apart from mathematics. Only with accurate measurement models can equity in learning analytics be realized. 

\section{Data and Code Availability}
The dataset used in the current study was collected under license of the third party company to which the data belongs. Restrictions apply to the availability of these data, but can be made available from the authors upon reasonable request and with permission of the company. Source code used to fit and analyze the data is publicly available at \url{https://anonymous.4open.science/r/qm-trees-10C5/}. 

%\bibliographystyle{apaparencite}
%\bibliography{qm_bib} 
\printbibliography

\end{document}
